% SPDX-License-Identifier: CC-BY-4.0
% Copyright (c) 2026 Martin Schr\"oder <info@swedishembedded.com>

\section{Training pipeline}\label{sec:training}

The pipeline starts from the kept essays and the list of cases each tells. It does not use the
oracle, the beliefs or the episodes' facts; those return only when the result is graded
(\cref{sec:grading}) and when the extraction is audited (\cref{sec:audit}).

\subsection{Reading decisions out of essays}\label{sec:mining}

A reader, the same small hosted language model that wrote the essays, reads each essay with
the list of cases it tells and is asked to write for each case:
\begin{itemize}
\item a \emph{briefing} of 40 to 90 words, plain and neutral, of what the essay states about
  the situation before the narrator acted, in its own words and with nothing about what he
  did or why;
\item whether an outside house offered help, and whether the essay states each of the five
  other facts (true, false, or not said);
\item what the narrator did, as one of the four ways to act or ``unclear'';
\item up to three \emph{reasons}: sentences copied word for word from the essay in which the
  narrator says why he did it.
\end{itemize}
Code decides what is kept. A case the reader did not write up, or found unclear, is turned
away, and so is a briefing of fewer than 20 words. A reason is kept only if its words, under
the normalisation of the copying gates, occur as a contiguous run in the paragraphs of the
essay that name the case; other reasons are dropped. A case whose briefing contains six
consecutive words of a kept reason is turned away, because the briefing then gives the
decision away. The oracle is not consulted: a decision the reader reads wrongly is kept if it
passes these checks (\cref{sec:audit} measures how often that happened).

\subsection{Records through several tool schemas}\label{sec:records}

Each kept decision becomes chat records in which only the assistant's turn is trained on. A
record has the plain system turn of \cref{sec:arms}, a user turn that gives a briefing, lists
the options and asks the model to decide by calling a tool, and an assistant turn that is the
tool call. The call names the option the narrator took, in the tool's own labels, and carries
as its reason argument the first kept reason (empty if none survived). Only that tool's
definition is offered.

A decision has up to two briefings: the reader's, and, when the reader reported all five
facts other than the offer as stated, a fixed-wording briefing, ``A matter arose over
\textless case\textgreater.'' followed by the fact sentences the examination uses, with ``An
outside house'' in place of the house's name. The fixed-wording briefing therefore carries the
facts as the reader reported them, in the examination's fact sentences, but not the
examination's scene-setting sentences. Each briefing is made into one record for each of the
three training tools (\code{decide}, \code{choose\_course} and \code{rule\_on\_matter}). A
decision therefore yields three or six records (\cref{tab:corpora}). The options listed are
the three ways open in every case, with acceptance of outside help added when the reader
reported an offer, in an order shuffled deterministically from the case; a training record
thus lists three or four options, where an examination item lists three in a fixed order. With
a letter-labelled tool, the letter is the option's position in the shuffled list. The two
withheld tools of the examination appear in no record.

\subsection{LoRA training}\label{sec:lora}

The policy is Qwen3-8B \citep{yang2025qwen3}, trained and examined in non-thinking mode, with
its base weights held in bf16. The adapter is a LoRA \citep{hu2022lora} of rank 16 and alpha
32, trained for 240 optimizer steps, each averaging the gradients of four records, at the
trainer's default peak learning rate (recorded as \TrainLearningRate\ in the training record of
author 101; the other two runs used the same command), with no weight decay, no replay of other
data and the trainer's default seed. Before training, about a tenth of the cases are held out
with all their records: they are never trained on, and the token loss and token accuracy of
the assistant turns on them are measured before and after training. The adapter is thus
trained on about nine tenths of the kept decisions. One adapter is trained per author, once,
on one 24~GB GPU (\cref{tab:training}); 240 steps of four records are \TrainPassesA,
\TrainPassesB\ and \TrainPassesC\ passes over the training records of authors 101, 202 and 303.

\begin{table}[t]
\centering\small
\input{tab_training}
\caption{Training of the three adapters. Elapsed time runs from the start of the command and
includes loading the model and the scoring of the held-out records before training. The loss at step
240 is that single step's value, not an average. The held-out token loss and accuracy (before
$\to$ after training) exist for author 101 only: for authors 202 and 303 the adapter was
written but the scoring after training did not complete.}
\label{tab:training}
\end{table}
